\section{Introduction}
\label{sec:introduction}

\subsection{Related Works}
\label{sec:related-work}

\textbf{Balancing scores and covariate balancing.} The propensity score is the canonical reduction of observed pretreatment covariates: under treatment ignorability given $X$, conditioning on $e(X)=P(A=1\mid X)$ balances the distribution of $X$ across treatment arms and identifies causal effects \citep{rosenbaum1983central}. Entropy balancing and the covariate-balancing propensity score instead choose weights or treatment models that target balance directly \citep{hainmueller2012entropy,imai2014covariate}. These methods motivate our optimization target, but their usual causal guarantee still requires the measured covariates to suffice for confounding control. Our setting keeps a latent confounder and asks what a held-out proxy view can certify about the bias remaining after a learned reduction of the other pretreatment measurements.

\textbf{Causal representation learning.} Representation-learning methods compress high-dimensional covariates while penalizing imbalance between treated and control representations. For example, counterfactual regression bounds counterfactual prediction error using factual prediction error and an integral probability metric between learned treatment-arm representations \citep{shalit2017estimating}. Our construction shares the use of a learned low-dimensional representation and a distributional balance penalty. The difference is the cross-view role of the proxy: for one prespecified block $j$, the learner constructs $Z_j=\phi_j(X,W_{-j})$ without outcomes and evaluates conditional treatment balance using the held-out variables $(X,W_j)$, even when ignorability given $X$ may fail.

\textbf{Negative controls and proximal causal inference.} Negative-control variables can reveal residual confounding because they should not be causally affected by the treatment or should not causally affect the outcome, while remaining associated with hidden common causes \citep{lipsitch2010negative}. Proximal causal inference turns this diagnostic idea into point identification by assigning distinct treatment-inducing and outcome-inducing proxy roles and imposing completeness and bridge conditions \citep{miao2018proxy,tchetgen2024ProximalCausalInference}; semiparametric proximal theory then develops influence functions and efficient estimators for the resulting identifying functionals \citep{cui2024semiparametric}. Our result has a different mechanism: a representation-side view forms the adjustment variable, a held-out view measures residual imbalance, and outcome-relevant completeness transfers exact observed balance to equality of the latent outcome averages needed for the population ATE.

\textbf{Single-proxy identification.} Single Proxy Control shows that one negative-control outcome can identify the effect of treatment on the treated under a counterfactual proxy condition and suitable bridge structure \citep{park2024single}. This differs from our population ATE target, from our latent-confounder proxy interpretation, and from our use of an observable balance discrepancy. The comparison is nevertheless central: it makes clear that ``one proxy'' does not by itself identify a common estimand. Identification depends on what the proxy measures, which conditional independence is imposed, and whether the analysis uses a bridge, an adjustment score, or an independent audit view.

\textbf{Kernel discrepancies and sensitivity certification.} Maximum mean discrepancy is an integral probability metric that equals zero for identical distributions when its kernel is characteristic \citep{gretton2012kernel}. We use conditional treatment-arm MMDs of $(X,W_j)$ as observable magnitudes, not kernel-test $p$-values. The new theoretical object is the product of this discrepancy and an outcome-relevant stability constant, which yields a bias certificate and a representation-selection guarantee. The individual ingredients have close precedents in balancing, representation learning, kernel testing, and proximal inference; the manuscript therefore makes no priority claim before a systematic literature audit, and instead states its contribution through the exact assumptions, estimand, and guarantee proved below.
